
\subsection{判別式に標識された五つの超平面での局所延長剛性}
\label{subsec:dp-eigen-hyperplane-rigidity}

四次del Pezzoの一般の$(S,[\sigma])$における全延長剛性は
定理\ref{thm:dp-rigidity}で既に示した。
ここでは、二次式pencilの判別式により自然に選ばれる
五つの特別な反標準切断について、\emph{一般性を仮定しない}
局所的な剛性を証明する。延長ファイバー全体の一意性とは区別する。

\begin{theorem}[判別式の固有超平面切断における被約孤立性]
\label{thm:dp-eigen-hyperplane-local-rigidity}
$S=V(q_0,q_1)\subset\PP^4$を任意の滑らかな
二つの二次式の完全交叉とする。
このpencilの判別式の五つの単純零点に対応する
五つの固有超平面を$H_0,\ldots,H_4$と書き、
$h_i\in H^0(\PP^4,\cO(1))$をその定義一次式とする。
各$C_i=S\cap H_i$は滑らかな楕円四次曲線である。
\[
 J_i=\tfrac14 J(q_0,q_1,h_i),\qquad
 \sigma_i=J_i|_S
\]
とする（$J$の正規化は式\eqref{eq:jacobian}とする）。

このとき任意の$i=0,\ldots,4$について、
正規化したPoisson延長スキーム
$\ExtP_S(\sigma_i)$の点$[J_i]$は
\emph{被約な孤立点}である。より強く、
\[
 \boxed{\quad T_{J_i}\ExtP_S(\sigma_i)=0,\qquad
       \cO_{\ExtP_S(\sigma_i),J_i}\simeq\CC.\quad}
\]
従って$J_i$の点は延長スキームの開かつ閉な連結成分をなし、
$\sigma_i$を固定したまま$J_i$から出発する非自明な
形式的Poisson延長変形は存在しない。

この結論だけでは、$J_i$から離れた別の成分に
階数四Poisson延長が存在する可能性を排除しない。
\end{theorem}
\begin{proof}
複素数体上、滑らかな二次式pencilを同時対角化すると
\[
 q_0=\sum_{j=0}^4x_j^2,\qquad
 q_1=\sum_{j=0}^4\lambda_jx_j^2,\qquad
 \lambda_j\ne\lambda_k\quad(j\ne k)
\]
と書ける。判別式の五つの零点はそれぞれ
$q_1-\lambda_iq_0$であり、固有超平面は$H_i=V(x_i)$。
その四次曲線$C_i$は$\PP^3$内で四つの相異なる
固有値の二つの対角二次式の完全交叉となるから滑らかである。
随伴によりその種数は一である。
$i=0$の場合を証明すれば十分である。
$q_1$を$q_1-\lambda_0q_0$へ置き換え、
$\lambda_0=0$としてよい。

$\mathcal E=\wedge^2T_{\PP^4}$とすると、
補題\ref{lem:kernel}により規格化した全双ベクトル延長は
\[
 P=J_0+q_0A+q_1B,\qquad
 A,B\in\wedge^2\CC^5
\]
で一意に表示される。$A=B=0$での
projective Jacobi方程式の微分を考える。
$J_0$は$q_0,q_1$をstrict Casimirとするため、
\[
 d\mathcal F_{J_0}(q_0A+q_1B)
 =q_0(E\wedge\ell_A(J_0))+
  q_1(E\wedge\ell_B(J_0)),
 \qquad \mathcal F(P)=E\wedge\Jac(P),
\]
となる。$\ell_A$は
式\eqref{eq:jac-factor}の直後で定義した定数双ベクトルとの
混合Schouten項である。
$q_0,q_1$は互いに素で、各
$U_A=E\wedge\ell_A(J_0)$は斉次二次の四重ベクトルである。
もし$q_0U_A+q_1U_B=0$なら
$U_A=q_1C,\ U_B=-q_0C$となる
定数四重ベクトル$C$がある。
$E\wedge U_A=0$から$E\wedge C=0$を得るが、
定数$C$でこの条件を満たすのは$C=0$だけである。
従って$U_A=U_B=0$である。

さらに$\ell_A(J_0)$は線形係数の三重ベクトルであり、
$J_0$と定数$A$は共に発散零なので
$\delta\ell_A(J_0)=0$である。
斉次座標のKoszul完全性から
$E\wedge\ell_A(J_0)=0$は
$\ell_A(J_0)=E\wedge T$という定数双ベクトル$T$を与える。
ところが五変数で
$\delta(E\wedge T)=3T$だから$T=0$であり、
$\ell_A(J_0)=0$となる。$B$についても同様である。

ここで$J_0$の定数双ベクトルによる混合Schouten項を
完全に計算する。$a<b$、$a,b\in\{1,2,3,4\}$を取り、
$\{c,d\}=\{1,2,3,4\}\setminus\{a,b\}$、$c<d$とすると
\[
 (J_0)^{ab}
  =\operatorname{sgn}(a,b,0,c,d)
       (\lambda_d-\lambda_c)x_cx_d,\qquad
 (J_0)^{0a}=0.
\]
定数双ベクトルの基底を
\[
 (01,02,03,04,12,13,14,23,24,34)
\]
の辞書式で取り、各$\ell_A(J_0)^{ijk}$の
一次単項式$x_m$の係数を行とする。
この$50\times10$行列で、
\[
\begin{array}{c|c}
\text{行番号（零始まり）}&(ijk;m)\\\hline
3,4&(012;3),(012;4)\\
9&(013;4)\\
19&(023;4)\\
31,32,33&(123;1),(123;2),(123;3)\\
36,37&(124;1),(124;2)\\
41&(134;1)
\end{array}
\]
を選んだ$10\times10$小行列式は正確に
\begin{equation}\label{eq:dp-eigen-nonzero-minor}
\boxed{\quad
-(\lambda_1-\lambda_2)(\lambda_1-\lambda_3)
(\lambda_1-\lambda_4)^2(\lambda_2-\lambda_3)
(\lambda_2-\lambda_4)^2(\lambda_3-\lambda_4)^3.
\quad}
\end{equation}
滑らかさから$\lambda_1,\ldots,\lambda_4$は相異なるので
これは零でない。
付属の独立producer
\texttt{check\_dp\_eigen\_rigidity.py}は
この普遍的な因数分解と、
$J_0$の発散・Casimir恒等式を
$\QQ(\lambda_1,\ldots,\lambda_4)$上で再計算する。
従って$\ell_A(J_0)=0$は$A=0$を強制し、
同様に$B=0$を強制する。
全20方向のJacobi線形化は単射である。

よって$\ExtP_S(\sigma_0)$の$J_0$における
Zariski接空間は零である。
有限型スキームのこの閉点の局所環はNakayamaの補題により
極大イデアルが零となり、$\CC$に等しい。
従ってこの点は被約かつ孤立し、
同じスキーム内で開かつ閉な連結成分となる。
他の$i$も座標の並べ替えで従う。
\end{proof}

\begin{remark}[特殊曲線を持つことと例外的延長は異なる]
滑らかな四次del Pezzoの判別式に由来する五つの
反標準楕円四次曲線は、標識された幾何学的に
特別な曲線である。しかし、これらに対応する
Jacobian型Poisson延長の局所スキームは全て被約一点である。
特殊な$h$であるという理由だけで、
余分なPoisson延長の分岐が生じるわけではない。
全延長スキームの他の孤立点や連結成分の不存在、
したがって普遍的な全点一意延長までは示していない。
\end{remark}
